\chapter{EDINET filings}\label{ch:edinet}

\section{Introduction}

EDINET is the electronic disclosure system of Japan's Financial Services Agency, the counterpart of the Commission's EDGAR, and it has required the financial statements of a securities report to be tagged in XBRL since 2008. Its API publishes three things: an index of what was filed on a given calendar day, the filing itself, and the agency's own conversion of the filing's XBRL into CSV \parencite{FSA.2026}. Module \emph{edinet} reads that conversion, just as module \emph{sec} of chapter \ref{ch:sec} reads the company files of the Commission and module \emph{esef} of chapter \ref{ch:esef} the xBRL-JSON of filings.xbrl.org. None of the three parses XBRL. The module builds three tables: one row per filing, one row per fact, and one row per filer of the agency's list of filers.

What arrives for a filing is a zip archive holding one CSV per document of it -- the report and the audit report, which the agency converts as well -- written in UTF-16 with a byte-order mark and separated by tabs, with nine columns and a header in Japanese. Unlike the company files of the Commission, the conversion carries the whole of a report: every fact, with its dimensions, its text, the Japanese label of its element, and the concepts the filer created for itself. Kojima's annual report for the year to 31 August 2026 comes to 1,099 facts, 1,086 of the report and 13 of the audit report, of which 19 are under concepts the filer created and 412 under the taxonomy of the financial statements. What the conversion leaves out is the presentation and calculation linkbases, an English label of anything but the filer's name, and -- the thing to know before using the tables -- the dates of a fact's context.

No column of the conversion holds the dates of a period. A context is named relative to the filing's own period, as in \verb|CurrentYearDuration| or \verb|Prior4YearDuration|, and a further column says the same in words. Two things turn such a label into dates, and the better one comes first. A filing states its own financial year and the comparative one among its facts, in the document and entity information, and those dates are used as reported, so the two periods nearly every use of the figures rests on are the filer's own dates rather than arithmetic. Only where no stated date covers the label -- the year before the comparative and everything further back -- is the period shifted by whole years, which is exact only where the filer's years are of equal length. A report stating five years of net sales is therefore anchored for two of them and shifted for three. A label whose shape is not recognised yields no dates at all rather than a guess, and the label is carried beside the dates so that what each rested on can always be read.

The index, the filings and the list of filers are downloaded when the pipeline runs and never shipped with the package. The documents are publications of the Financial Services Agency, whose reuse the agency permits under the terms it publishes on the disclosure site. The API refuses a request that carries no registered key. The key is read from an argument, an environment variable or a file in a directory of secrets, and never from the configuration file of the data lake: that file is committed with a deployment, handed to every module and opened by every report template, so a key in it would travel with the repository. The key travels as a query parameter, so the address of every request contains it, and nothing in the module logs or records an address before the key has been replaced in it. The list of filers is the exception: it is a zip on the disclosure site and needs no key, which is why a lake can name its filers before it can collect a filing from any of them.

\section{Data pipeline}

\subsection{Landing zone}

There is no landing zone. An index is a page of JSON and an archive a few hundred kilobytes, both written into the Bronze layer once they are known to be what was asked for.

\subsection{Bronze layer}

The Bronze layer differs in kind from those of the other two regimes, and the reason is that a filing leaves EDINET. A securities report is on public view for five years and retrievable for five more, and the specification says plainly that what is older cannot be obtained. No parameter of the API reaches past that, and a request for a day outside the window returns an empty index rather than an error, which is how a lake that believes it is collecting ten years of history quietly collects what is left of it. EDGAR keeps a filing indefinitely and filings.xbrl.org keeps its archive; EDINET does not. The Bronze layer is therefore not a convenience that saves a download but the only lasting copy, and a year not collected is a year lost rather than a year postponed. The pipeline refuses a day or a filing the window no longer covers, with a message that says so, rather than storing the empty answer the API gives for one.

The index is the unit of collection. The API answers for one calendar day, not for one filer and not for a range, so a lake that wants a year asks 365 times and a lake that wants the retrievable history some 3,600 times, plus one request per filing converted. Every index fetched is kept, so a second run over a year the lake already holds makes no request at all. A run brings the last thirty days up to date by default, which is what keeping a lake current costs; a run that wants the history passes the earliest day still retrievable.

The API reports a failure in two ways, and neither is the HTTP status. The index API answers an error with status 200 and the real status in the body: a request with no key, with an unknown key and with an inactive one all come back as 200 carrying a 401 inside. A run that trusted the status would store that as the index of the day and report success, so the body is read and the run raises on what it finds there. The document API answers an error with a different content type, an archive for a filing and JSON for a failure, which the specification names as the way to tell them apart. A body that is not an archive is refused rather than stored, since an error page written into the Bronze layer would otherwise be the only copy of a filing that no longer exists.

\subsection{Silver layer}

The table of filings has one row per document the collected indexes name, partitioned by the year it was filed, and is rebuilt in full from every index the Bronze layer holds. A filing is identified by the document management number the API addresses it by. The index says per filing whether the agency converted its XBRL, and a filing it did not convert contributes nothing to the table of facts, exactly as an ESEF filing without xBRL-JSON. The document type says what kind of document it is, an annual securities report being type 120; most documents of a day are not company reports at all but large-holding reports, fund reports and the like. Four states travel with a filing and are carried rather than acted on: withdrawn, corrected by a finance bureau, withheld from disclosure, and whether it is on public view, in the extension period or past both. A filing is one row however often it is listed. The agency lists a filing again, on the day of the change, whenever its details are edited, and the table keeps the last listing, which is the filing's present account; keeping every listing would read the archive of such a filing once per listing and multiply its facts in every table built on this one.

The table of facts has one row per fact of a converted filing, partitioned by the year the filing was made. That year only moves forward, as a financial year does not, since one report states five of them, so a run that collects this month's filings rewrites this year's partition and leaves the others alone. The archives are read on the executors, one at a time, and an archive that cannot be read comes back as one row carrying its traceback, which the driver logs and counts. A filing listed with a conversion but not collected is counted and reported, since that is a gap a run can close. The columns a fact shares with the other two regimes are spelled as they spell them. Three columns have no counterpart there: the Japanese label of the element, the dimensions of the context, and the basis, which says whether a figure is the group's or the parent company's. A fact can have a unit and no number, the conversion writing a mark for not applicable where a figure does not apply, and that becomes an empty value rather than a zero. A fact is identified within its filing by the document, the concept, the context and its occurrence; a fund manager files one document per series of a fund under one form, with different values for the same element, and the series is then part of the identifier.

A Japanese report states its figures twice, and no reader should take the first row that matches. The consolidated accounts carry no dimension and the parent company's own carry a member for the non-consolidated basis; Kanagawa Chuo Kotsu's annual report states earnings per share of 294.99 for the group and 236.58 for the parent in the same table of the same report. A filer that prepares no consolidated statements states its accounts under that member alone, so a query that demands an empty dimension returns nothing for it. Whether a filer consolidates is stated in the filing itself, and the agency's list of filers carries a flag of its own that can disagree with it.

The table of filers has one row per filer the agency lists, the counterpart of the Commission's list of registrants and the way a lake finds the codes of the companies it wants by name, ticker or industry. In November 2026 it held 11,467 filers, 3,816 of them listed. The financial year end is carried as text, as the agency writes it, because there is no year in it and because the agency writes the end of March in two ways. It is worth having: 1,633 of the 3,816 listed filers do not close in March, and a comparison that assumes the Japanese financial year ends there is wrong for every one of them. The list also carries the corporate number of the National Tax Agency, which section \ref{sec:edinet-lei} returns to.

\subsection{Data quality}

Twenty-six checks write what they found into a fourth table of the Silver layer, \emph{issues}, with the machinery shared with the other two regimes. As there, nothing is dropped and nothing is corrected. The arithmetic checks are the valuable ones, and two of them exist only here, because they are the two lines that make a Japanese profit and loss account what it is: the ordinary result is the operating result and the recurring items outside it, and the result before tax is the ordinary result and the extraordinary items. Neither relation can be written for an IFRS or a US GAAP filer, which has neither line.

Four further checks have no counterpart anywhere in the package, and each follows from a property of the source described above. The five-year summary of the front matter and the financial statements both carry the revenue and the total assets of the year the report is about, in different elements of different taxonomies, so they have to agree; this is the only check in the package that tests a filing against another part of itself. The years older than the comparative are shifted by whole years, and since the filing states the first and last day of its two most recent years, whether the shift can be trusted is a question the filing answers, with three days of slack so that a leap year is not taken for a 52-week calendar. The basis is stated twice, as a column and as a dimension, and the two must not disagree. And an element named for a ratio must be a pure number: this is the one check of the four that has found anything in a real report, Fujitsu's, which tags its equity per share with 1,167.25 under the element of the equity ratio, so that a reader trusting the name would print an equity ratio of 1,167 per cent. The ratio itself is in the same report under another element as 0.60.

Two decisions are settled once rather than by each check. The figures of a filing are taken on the basis the filing itself says it reports on, consolidated where it consolidates and the parent's own where it does not, and a fact carrying a further member is a breakdown and left out. And a missing line is not a zero: a component of a sum being tested has to be stated, and a difference that cannot be computed is not a difference. Treating an absent line as nothing turns a filing that presents its statements differently into one that fails arithmetic; the first draft of the checks reported Fujitsu's liabilities as short by 244,695 million because its IFRS statement leaves the non-current liabilities to be read as the remainder. Only a genuinely optional line, such as extraordinary items a filer may not have had, is read as nothing. Over four real annual securities reports -- Kojima, which does not consolidate; Kanagawa Chuo Kotsu twice, consolidated under Japan GAAP; and Fujitsu, consolidated under IFRS -- the checks report five issues, all five Fujitsu's mis-tagged ratio.

\subsection{Gold layer}

There is no Gold layer. An \emph{analyst factsheet} on one filer reads the Silver tables as Parquet, as its ESEF and SEC counterparts do, and does what a Gold table would do for its question. It keeps the two spans of a Japanese report apart: the five-year summary of the front matter gives five years of headline figures from a single report, which neither of the other regimes affords, while the statements give the year and its comparative. A concept is matched on its local name and never on its prefix, since the fourth prefix of a conversion is built from the filer's own code and carries elements the core taxonomy lacks, such as an IFRS filer's operating profit. Every line declares the class of unit it accepts and refuses a fact of another, which is what keeps Fujitsu's equity per share out of its equity ratio. No figure per share is adjusted for a split, because a Japanese filer restates the earlier years of its own summary at its discretion and not always for every year; the year the share count moved is named instead.

\section{Software architecture}

Module \emph{edinet} sits in the package \emph{financialreport} and is organised like the other two. The part \emph{common} declares the layout of the data lake and the schemas of the tables; \emph{parse} reads an index into filings and an archive into facts in plain Python, resolving the periods as described above; \emph{bronze} collects the list of filers, the indexes and the archives; \emph{silver} builds the three tables; \emph{quality} holds the checks; and \emph{pipeline} runs the whole in one call,
\begin{quote}
\verb|edinet.update_edinet(spark, "paths.yaml")|
\end{quote}
which collects the indexes of the last thirty days, then the conversions of the filings those indexes say were converted, then rebuilds the tables. A filer is keyed by its EDINET code, the letter E and five digits. The steps can also be run on their own, and the checks are a call of their own, \verb|update_silver_edinet_quality|.

\subsection{The corporate number reaches the LEI}\label{sec:edinet-lei}

Of the three regimes, only the Japanese one can be joined to the LEI, and the reason is that the corporate number of the National Tax Agency is not an identifier of its own. It is a check digit followed by the twelve-digit number of the commercial register, and GLEIF records that register number as the registration authority entity identifier of an LEI issued to a Japanese company. The join is therefore arithmetic rather than a lookup table: drop the check digit, hyphenate the rest as the register writes it, and look it up in the entity master data of chapter \ref{ch:entitymasterdata} under the register of the Legal Affairs Bureau. Because it compares two spellings of one register number and not names, it survives a rename: of 27 matches found in a sample, one filer had renamed itself since its LEI was issued, and the number matched where the name would not have.

It reaches far less than it promises. Of a random sample of 200 listed filers with a corporate number, 27 reached an LEI, and ten of the misses had no LEI under any spelling of their name: most Japanese listed companies are not in the LEI system at all. The join is enough to put a Japanese filer beside a European one where both exist, which the combined tables of chapter \ref{ch:financialreport} cannot otherwise do, and not enough to key anything by. The EDINET code stays the key for Japan, as the central index key does for the Commission, and nothing in the module derives an LEI.
